\section{Related Work}
\label{sec:related}

\subsection{The Missing Join}
The closest systems do observe many of the ingredients of governance drift,
but they organize them around plane-local questions: desired versus observed
configuration, state versus current policy, artifact versus admission rule,
or environment versus benchmark. \Cref{tab:scope-comparison} makes the gap
explicit. The table compares each line of work by its \emph{native decision
object}, not by every integration a commercial product might expose. A
triangle means that the line carries a useful slice but does not natively join
that slice to the time-indexed admitted basis.

\begin{table*}[t]
\centering
\caption{Representative Native Decision Scope of Adjacent Work. \cmark{} =
Native Decision Object; $\triangle$ = Partial Input or Slice; -- = Not Native.
The Comparison Is Not an Exhaustive Product Feature Audit.}
\label{tab:scope-comparison}
\scriptsize
\setlength{\tabcolsep}{3.2pt}
\begin{tabular}{@{}lccccccc@{}}
\toprule
Line of work & Config. & Policy & Auth. & Intent & Evidence & Environ. & Joins admitted basis \\
\midrule
Desired-state/GitOps reconciliation~\cite{argocd,fluxcd,opengitops} & \cmark & -- & -- & -- & -- & -- & -- \\
Security-focused configuration baselines~\cite{nist800128,nist80053} & \cmark & $\triangle$ & -- & $\triangle$ & $\triangle$ & $\triangle$ & -- \\
IaC drift/configuration recorders~\cite{awsconfig,driftctl} & \cmark & -- & -- & -- & -- & $\triangle$ & -- \\
AWS Control Tower governance drift~\cite{awsctgovernancedrift} & \cmark & \cmark & $\triangle$ & -- & -- & $\triangle$ & -- \\
Multi-cloud policy consistency/drift~\cite{marella2026multicloud} & \cmark & \cmark & -- & -- & -- & \cmark & -- \\
Policy-as-code/background scans~\cite{opa,kyverno} & -- & \cmark & -- & -- & -- & -- & -- \\
Posture and hardening scans~\cite{cspm,nsacisa2022} & $\triangle$ & \cmark & -- & -- & -- & \cmark & -- \\
Supply-chain attestations~\cite{slsa,newman2022} & -- & $\triangle$ & $\triangle$ & $\triangle$ & \cmark & -- & -- \\
Continuous artifact validation~\cite{binauthz} & -- & \cmark & \cmark & -- & $\triangle$ & -- & -- \\
\textbf{Governance drift (this work)} & \cmark & \cmark & \cmark & \cmark & \cmark & \cmark & \cmark \\
\bottomrule
\end{tabular}
\end{table*}

This fragmentation explains why the aggregate relation remained easy to
miss: a deployment can be green in each available local view while the join
among present state, current context, and the recorded admitted basis is
broken. The contribution is not the claim that no existing mechanism can be
composed to cover several columns. It is the explicit temporal three-way
relation, the class-resolved verdicts it yields, and the evidence coverage
required to decide it.

\subsection{Terminology and Review Boundary}
The phrase \emph{governance drift} predates this work, and we make no claim to
have coined it. AWS Control Tower officially uses the phrase (also
``organizational drift'') for changes to managed organizational units,
accounts, landing-zone artifacts, controls, policies, and inherited
baselines~\cite{awsctgovernancedrift}. That product taxonomy is an important
precedent. Its documented decision object is nevertheless Control
Tower-managed organizational and control state; AWS explicitly excludes
resource drift inside a baseline. It does not define a deployment-specific
join between present workload state, evolving governance context, and a
uniquely activated immutable approval record.

Two 2026 papers further sharpen the boundary. Marella's SSRN preprint proposes
centralized policy abstraction, cross-cloud policy mapping, and continuous
drift detection across AWS, Azure, and GCP~\cite{marella2026multicloud}. Its
object is provider-independent policy consistency; that layer can supply T1
and T4 evidence, but does not by itself establish whether a running deployment
remains covered by the authorization, intent, and environmental assumptions
recorded at admission. Do et al. use the exact phrase \emph{AI governance
drift} for post-deployment misalignment between governance mechanisms and the
behavior of deployed AI systems, and evaluate an organizational model relating
audit, monitoring, assurance, visibility, drift, and resilience
\cite{do2026aigovernancedrift}. Their unit of analysis is AI governance
capability and behavior rather than a cloud deployment's immutable admitted
basis; for that reason, the work is discussed here but is not forced into the
mechanism columns of \cref{tab:scope-comparison}.

Other 2026 uses confirm that the compound is polysemous rather than available
for a naming claim. H\"opfner and Schacht study \emph{socio-technical
governance drift} as longitudinal erosion of rule compliance in human--AI
teams~\cite{hoepfner2026architecture}. Solozobov's preprint and accompanying
toolkit monitor degradation of the statistical evidence supporting deployed
risk models under delayed ground truth~\cite{solozobov2026labelfree}. The
former concerns organizational oversight behavior; the latter concerns model
performance evidence. Neither makes cloud workload admission history its
decision object, but both are relevant neighbors and are included rather than
dismissed as mere keyword collisions.

We tested the remaining novelty statement with a bounded, one-reviewer
\emph{scoping review} updated through August 8, 2026---not a systematic review
or a claim of corpus completeness. The protocol records the exact queries,
sources, cutoff, inclusion criteria, and limitations; its item-level ledger
records the included and excluded candidates and a reason for every decision
(\path{literature/scoping_review_protocol.md} and
\path{literature/scoping_review_records.csv}). Discovery used exact-term
and adjacent-construct queries, while substantive claims were checked against
official documentation, standards, DOI/publisher records, and primary
proceedings. No corpus-size estimate, dual screening, or risk-of-bias appraisal
is implied. The bounded conclusion is correspondingly narrow: among the
screened sources, none jointly specifies the full deployment-level,
history-anchored, class-resolved relation defined here.

\subsection{Configuration Drift and Desired-State Systems}
Configuration drift detection is mature practice: reconciler status in
GitOps~\cite{argocd,fluxcd,opengitops,beetz2022}, cloud configuration
recorders~\cite{awsconfig}, IaC drift tools~\cite{driftctl}, and the
configuration-management control tradition~\cite{nist800128,nist80053} that
institutionalizes baseline comparison. IaC research studies the quality and
smells of the configuration artifacts themselves~\cite{rahman2019mapping,
rahman2019,sharma2016}. GitOps reconcilers and many deployed drift tools use
the current desired state as their comparator, but configuration management as
a field is not uniformly memoryless. NIST SP 800-128 requires an established
baseline, security-impact analysis, approval and documentation of controlled
changes, and monitoring of the resulting configuration~\cite{nist800128}; it
therefore supplies genuine approved-history semantics. The narrower boundary
is dimensional: that history is configuration-scoped, whereas our relation
joins the selected deployment admission to current policy, authorization,
artifact evidence, intent, and environment as well. \Cref{prop:impossible}
bounds what the configuration-only vocabulary can decide, and
\cref{fnd:blind} measures the bound. Our claim is containment, not erasure:
this literature solves the configuration component well---T0n in our
architecture \emph{is} this literature---but it does not decide all six
classes from configuration evidence alone.

\subsection{Policy-as-Code, Admission, and Continuous Compliance}
Policy engines~\cite{opa,kyverno} and admission
control~\cite{k8sadmission} evaluate governance predicates at a point in
time; posture management~\cite{cspm} and hardening baselines~\cite{nsacisa2022}
scan environments against benchmark policy; continuous-validation features in
attestation-gated deployment~\cite{binauthz} re-check the artifact-policy
slice after admission. These supply, respectively, our T1 evaluation
semantics, part of T4's observation, and a one-tier precedent for continuous
re-evaluation. Commercial continuous-compliance and CSPM products may
integrate several of these signals, but their native decision object remains
current state versus current rule or benchmark; this is not an exhaustive
feature audit. What the documented mechanisms do not provide is the
\emph{approved-state reference}:
they compare state against current policy, not against the basis under which
the state was admitted, so they cannot distinguish ``was never compliant''
from ``was admitted legitimately and the world moved''---the distinction
carrying the remediation semantics of \cref{sec:model-vector}, and the
defining relation of governance drift.
Runtime compliance-monitoring research systematically studies rule
evaluation over process-event streams and fine-grained violation feedback
\cite{ly2015compliance}. That literature supplies monitoring languages and
execution semantics; its native reference is a compliance rule or process
model, not a deployment-specific immutable admission event. The distinction
is complementary rather than competitive: such monitors can implement T1,
while $B_x(t)$ determines why the same current violation is classified as
post-admission drift rather than absence of an admitted basis.
Within the screened set, the closest new result on the authorization timeline
is Peng and Wu's policy-state serializability: it prevents an authorization or
delayed human approval from becoming stale while mutable policy state changes
before the governed effect commits~\cite{peng2026stateful}. Its protected interval ends at
effect commit; ours begins from the admitted deployment and asks whether the
running object remains covered as multiple governance planes evolve afterward.
The guarantees are complementary: serializable admission can create a sound
basis, and the retained basis can support post-admission drift detection.

\subsection{Supply-Chain Integrity and Authorization Evidence}
Attestation frameworks bind artifacts to build and review
facts~\cite{slsa,torres2019,newman2022}; they are T3's substrate and make
S4/S8-class drift decidable. Their orientation is admission-time
verification; governance drift adds the post-admission temporal dimension
(validity windows, revocation, supersession) and the join against
$\Gapp$. Howell and Kotz's end-to-end authorization carries the authority
justifying a request across administrative, network, abstraction, and protocol
boundaries~\cite{howell2000e2e}; governance drift preserves an analogous
justification across time, but asks whether a previously admitted, still
running deployment remains covered after its context changes. Provenance work
provides equally direct foundations for the evidence join. Muniswamy-Reddy
et al. integrate provenance across application, workflow, storage, and network
layers~\cite{muniswamyreddy2009layering}; TAP explicitly represents time,
distributed state, and state changes in provenance~\cite{zhou2011tap}.
Those systems solve cross-layer ancestry and temporal explanation, not the
normative, class-resolved conformance test against an activated approval
record. Records-side, durable approval evidence and transparency
services~\cite{scitt2025,provdm,cloudtrail} determine whether that test is
decidable at all---our observability mask and GCC metric make the dependence
explicit rather than assumed.

\subsection{Access Control, Runtime Verification, and Drift Theory}
Temporal validity of authorization has deep roots in temporal authorization
and role models~\cite{bertino1994temporal,bertino2001trbac}, while usage
control adds ongoing decisions, mutable attributes, obligations, and
conditions~\cite{park2004ucon,schoepp2023usage}. These works justify the
one-shot/continuing distinction rather than leaving validity as a Boolean.
Classical access-control theory~\cite{sandhu1996,clark1987} and zero-trust's
continuous-evaluation stance~\cite{nist800207} provide the broader foundation;
governance drift applies it to deployment admission history. Runtime
verification~\cite{leucker2009,pnueli1977} monitors executions against
temporal specifications; our evaluator is kin---governance consistency is a
safety-style property over joined state and context streams---but the
specification here is not a given formula: it is reconstructed per deployment
from the approved snapshot, which is why recording $\Gapp$ is the
architecture's prerequisite rather than an optimization. Concept
drift~\cite{gama2014} shares only the word: it names distribution change in
data streams, not divergence from an approval basis; we flag it to prevent
terminological import of its detectors, which answer a statistical question
ours does not ask. The exact-phrase precedents and review boundary are stated
above so that the contribution rests neither on naming priority nor on a
universal absence claim. It rests on the specified decision object and the
analytic consequences of retaining the admitted basis.

\subsection{Novelty, Answered Directly}
\emph{Is this just configuration drift?} No---with the argument's weight
placed where it belongs. The load-bearing evidence is analytic: constructed,
tool-realizable counterexamples in which configuration equality and
governance inconsistency coexist (\cref{fig:example}), and the dependency
result that no configuration-pair detector decides the property
(\cref{prop:impossible}). The executed staircase (\cref{tab:matrix})
supports these as an implementation-validation and noise study, not as
independent field evidence. And we concede the mechanical point a skeptic
will press: every governance-tier detector is itself a comparison---against
a recorded basis, a validity clock, or a coverage set. The claim is not a
new mechanism but a different \emph{relation}: three-place
(state, current context, admitted basis) where the desired/observed
configuration signal is two-place and memoryless, with the consequences \cref{sec:discussion}
enumerates. \emph{Is it policy compliance scanning?} Scanning
evaluates current state against current policy; governance drift is
three-way---state, current context, \emph{and approved basis}---and the third
term changes both the verdicts (legitimately-admitted-since-superseded
$\ne$ never-compliant) and the remediations. \emph{Is it continuous compliance scanning or continuous validation?}
Per tier, largely yes---and we credit it: Kyverno-style background scanning
continuously re-evaluates resources against current policy (T1);
configuration recorders and posture management watch unmanaged resources
(much of T4); continuous validation re-checks artifact policy (part of
T3)~\cite{kyverno,awsconfig,cspm,binauthz}. Several of our scenarios
trigger \emph{some} existing detector. Within the bounded scoping set, we
found no documented mechanism providing the complete join against the
admitted basis with mode-aware validity and
class-resolved verdicts---the increment that turns ``a scanner fired'' into
``this running state is no longer covered by what admitted it, for this
reason, with this remediation family.'' That increment is architectural and,
we argue, decision-relevant; quantifying its operational value on real
estates remains future field work. The present laboratory tests only the
three hypotheses stated in \cref{sec:lab}.
